\documentclass[11pt]{amsart}
\usepackage[a4paper,margin=1in]{geometry}
\usepackage[T1]{fontenc}
\usepackage{lmodern,microtype,amsmath,amssymb,booktabs,array}
\usepackage[hidelinks]{hyperref}
\newtheorem{theorem}{Theorem}[section]
\newtheorem{lemma}[theorem]{Lemma}
\newtheorem{proposition}[theorem]{Proposition}
\theoremstyle{remark}\newtheorem{remark}[theorem]{Remark}
\newcommand{\R}{\mathbb R}
\newcommand{\diag}{\operatorname{diag}}
\newcommand{\tr}{\operatorname{tr}}
\newcommand{\one}{\mathbf 1}
\newcommand{\eps}{\varepsilon}
\newcommand{\T}{\mathsf T}
\newcommand{\calU}{\mathcal U}
\newcommand{\calD}{\mathcal D}
\title[Constructive Koml\'os discrepancy below 67]{Constructive Koml\'os discrepancy below 67:\newline coupled spectral signing with input-linear row processing}
\author{}
\date{27 September 2026; revised complexity analysis}
\begin{document}
\begin{abstract}
We modify the deterministic retained-row spectral walk of the project
draft~\cite{project}. The discrepancy bound becomes
$1271/25+(149/100)\sqrt{117}=66.95681420\ldots<67$.
The changes are a coupled Gram--resolvent estimate, a nonquadratic concave
coordinate profile, and polynomial-time terminal rounding that uses each
row's unspent deletion allowance. The profile minimizes its height subject
to the particular pointwise gradient--curvature condition used in this
proof. We extend the finite Perron comparison to its $C^{2,1}$ weight
curves and retain $O(n\log^4(2+n))$ updates.
Combining this walk with Li's row screening and motion clock gives
$O(mn)+\widetilde O(n^4)$ deterministic field operations.
The sharper fast-arithmetic bound is
$O(mn)+\widetilde O(m_0n^2+n^{\omega+1})$, where
$m_0< n^2/66^2$ is the number of rows surviving a single input scan.
The fast version uses the draft's imported linear-algebra primitives and
square roots. No polynomial bit-complexity or practical runtime claim is made.
\end{abstract}
\maketitle

\section{Result and the source of the improvement}

\begin{theorem}\label{thm:main}
Let $A\in\R^{m\times n}$, $n\ge1$, satisfy $\sum_i a_{ij}^2\le1$ for every column.
There is a deterministic algorithm that returns one sign per original
column, $\eps\in\{-1,1\}^n$, with
\[
 \|A\eps\|_\infty<C_*:=\frac{1271}{25}+\frac{149}{100}\sqrt{117}
 =66.956814201324\ldots<67.
\]
Let $m_0=|\{i:\sum_j|a_{ij}|>66\}|<n^2/66^2$.
Put $J=1+\lceil\log_2 n\rceil$ and
$q=\lceil\log_2(8(n+2m_0))\rceil$ for $n\ge1$.
The algorithm uses $O(nJ^2q^2)=O(n\log^4(2+n))$ updates and
$O(mn)+\widetilde O(n^4)$ field operations and comparisons.
Its states are rational for rational input. With the fast symmetric
diagonalization and exact rank primitives specified in the source draft,
and allowing square roots, its cost is
\begin{equation}\label{eq:screenedmaincost}
 O(mn)+\widetilde O(m_0n^2+n^{\omega+1}).
\end{equation}
All logarithms hidden by $\widetilde O$ depend only on $2+n$.
\end{theorem}

The underlying retained-row walk, affine offsets, row deletion, restricted
trace argument, exact constraint projection, deterministic flattening,
and phase regularization are those of the project draft~\cite{project},
whose deterministic core is also recorded in~\cite{core}. The coupled
Gram--resolvent inequality used below appears as Lemma 2.4 of Li~\cite{Li};
we include its elementary proof. The new ingredients here are its joint
use with the profile in Section~\ref{sec:profile}, the remaining-mass
rounding argument, and the checks that make the altered walk finite.
We make no priority claim for the elementary spectral inequality.

The complexity refinement in Section~\ref{sec:screening} borrows Li's
warm/cold screening, motion clock, and aggregate statistics, including
their elementary identities. In place of his quadratic potential charge,
we control row energy directly by two additional direction constraints.
The further work is to fit the screening into the below-$67$ walk without
spending its discrepancy allowance, handle unflattened cold derivatives,
and bound the clock using the nearly linear update count. The resulting support
maintenance cost is $O(m_0n^2)$ and the remaining fast work is
$\widetilde O(n^{\omega+1})$.

Sections~\ref{sec:profile}--\ref{sec:arithmetic} give the unscreened core
and its correctness proof. Section~\ref{sec:screening} specifies the
screened implementation proving the improved running time in
Theorem~\ref{thm:main}. Preprocessing discards rows of $\ell_1$ norm at
most 66. In all walk formulas below, row indices refer only to the surviving
matrix; any row-count variable $m$ in a core cost is consequently $m_0$.

\begin{center}
\begin{tabular}{lrr}
\toprule
Proof contribution & Previous deterministic core & This supplement\\
\midrule
Tracked-row height & $75$ & $34.45$\\
Total deletion allowance & $59.849623\ldots$ & $32.506814\ldots$\\
Resulting bound & $134.849623\ldots<135$ & $66.956814\ldots<67$\\
Terminal rule & nearest signs, $s\le99$ & conditional expectation, $s\le512$\\
\bottomrule
\end{tabular}
\end{center}

The separate cosine-translation overlap theorem in the project's existence
work does not by itself give polynomial-time access to signing directions.
There is no such transfer in this proof. Here a sharp local
gradient--curvature relation makes the directions computable within the
spectral walk. The theorem is an arithmetic-operation statement; it does
not establish bounded bit lengths or an efficient numerical implementation.

\section{A profile matched to the curvature estimate}\label{sec:profile}

Set
\[
 d=\frac2{11},\qquad c=\frac{17}{25},\qquad
 f(x)=\frac{275}{34}\log\frac{25-17|x|}{8}
                 +\frac{11}{2}(|x|-1),\quad -1\le x\le1.
\]
The function is even, nonnegative and concave, vanishes at $\pm1$, and is
$C^2$ with Lipschitz second derivative. In particular,
\begin{align}\label{eq:fderivatives}
 f'(x)&=-\frac{187x}{2(25-17|x|)},&
 -f''(x)&=\frac{4675}{2(25-17|x|)^2},\notag\\
 \|f'\|_\infty&=\frac{187}{16},&
 \|f''\|_\infty&=\frac{4675}{128},&
 \|f'''\|_{\infty,\mathrm{a.e.}}&=\frac{79475}{512}.
\end{align}
At zero the formula for $f''$ is interpreted continuously. Its maximum is
\[
 f(0)=3.716012584611\ldots<a_0:=\frac{93}{25}.
\]
Most importantly,
\begin{equation}\label{eq:exactratio}
 \frac{(1+d|f'(x)|)^2}{-f''(x)}=\frac dc=\frac{50}{187}.
\end{equation}
Consequently, if $z=t^2\le d^2$, then for either sign $\sigma$,
\begin{equation}\label{eq:pointwise}
 (\sigma t+z f'(x))^2
 \le z(1+d|f'(x)|)^2
 =\frac{50}{187}\,z(-f''(x)).
\end{equation}

\begin{proposition}[Optimality within the pointwise profile condition]
Fix $d>0$ and $\lambda>d$. Among even, differentiable concave functions
$F$ with $F(1)=0$ satisfying
$(1+d|F'|)^2\le\lambda(-F'')$, the smallest possible value of $F(0)$ is
attained by
\[
 F(x)=\frac1d\left[\frac1c\log\frac{1-c|x|}{1-c}+|x|-1\right],
 \qquad c=d/\lambda.
\]
The derivatives may be understood almost everywhere, provided $F'$ is
locally absolutely continuous. This is optimality for this specified
condition, not a global optimum over signing algorithms.
\end{proposition}
\begin{proof}
On $[0,1]$ put $g=-F'\ge0$, so $g(0)=0$. The hypothesis gives
$g'\ge(1+dg)^2/\lambda$. Hence
$(1+dg(x))^{-1}\le1-dx/\lambda$, and
$g(x)\ge x/(\lambda-dx)$.
Integrating $F(0)=\int_0^1g(x)\,dx$ gives the displayed profile's height.
Equality solves this differential equation. A finite height is impossible
under this condition when $\lambda\le d$.
\end{proof}

\section{State, parameters, and cleanup}\label{sec:state}

The constants used throughout are
\begin{equation}\label{eq:parameters}
\begin{gathered}
 R=\frac{149}{100},\quad \Delta=d^2=\frac4{121},\quad
 \beta=\frac45,\quad a_0=\frac{93}{25},\quad b=\frac{218}{25},\quad
 H_0=\frac{689}{20},\\
 A_0=36,\quad B_0=\frac4{625},\quad
 \theta=\frac{1001}{1000},\quad L_* =\frac{257}{256},\quad
 \kappa=\frac{47}{200},\quad \alpha=\frac{59}{250}.
\end{gathered}
\end{equation}
Start at $x=0\in[-1,1]^n$. Let $V$ be the active coordinates, $s=|V|$.
Row $i$ retains a set $T_i\subseteq V$ and an offset $c_i$; initially
$T_i=V$ and $c_i=0$. Write
\[
 b_{ij}=a_{ij}\one_{j\in T_i},\quad p_{ij}=b_{ij}^2,\quad
 q_i=\sum_jp_{ij},\quad d_i=c_i+b_i^\T x_V,\quad
 E_i=\sum_jp_{ij}f(x_j).
\]
Removing a retained coefficient transfers $a_{ij}x_j$ to $c_i$, preserving
$d_i$. We always have $0\le E_i\le a_0q_i$ and $\sum_iq_i\le s$.
A row is large if $q_i>R^2$, medium if $0<q_i\le R^2$, and completed if
$q_i=0$. Large rows impose $b_i^\T h=0$ and undergo no deliberate deletion.
On medium rows delete entries with $p_{ij}>\Delta q_i$, one at a time,
updating $q_i$ and choosing the least eligible pair. Thus a cleaned medium
row has $\max_jp_{ij}\le\Delta q_i$.

A nonempty medium row has a nonincreasing rational scale with
$q_i\le r_i^2<\theta^2q_i$ and $r_i\le R$. Compute a candidate
$\theta^\ell\max_{j\in T_i}|a_{ij}|$, using the least nonnegative integer
$\ell$ whose square is at least $q_i$, and take the minimum of this
candidate, $R$, and the previous scale if present. This requires
$O(\log(2+n))$ operations, with an absolute constant depending on $\theta$.

For medium row-sign pairs set
\begin{equation}\label{eq:barrier}
 u_{i\sigma}=\frac{\sigma d_i-H_0}{r_i}+\frac{E_i}{r_i^2}+b,
 \qquad w_{i\sigma}=A_0r_i^{-4}e^{\beta u_{i\sigma}}.
\end{equation}
The invariants are $u_{i\sigma}\le0$, $d_i=0$ on large rows, and
$L<L_*$. Here
\begin{equation}\label{eq:matrix}
 W=\sum_{i,\sigma}w_{i\sigma}p_ip_i^\T,\quad
 B=B_0\sum_{i\text{ large}}\diag(p_i),\quad
 M=W+B+\eta\one\one^\T,\quad L=\lambda_{\max}(M).
\end{equation}
For an empty active set set $L=0$. During a phase the potential is
\[
 \Phi=(L-1)_+^3-\eta\|x\|_2^2,
\]
where the norm includes frozen coordinates.

Cleanup freezes boundary coordinates, absorbs their retained contributions,
introduces newly medium rows, performs deliberate medium deletions, and
decreases scales. Row classes can only decrease.

\begin{lemma}\label{lem:cleanup}
Cleanup preserves tracked sums, does not increase $L$ or $\Phi$, and
preserves feasible barriers. Initialization has
$L\le B_0+\eta n<1$ when $\eta n\le2^{-40}$.
\end{lemma}
\begin{proof}
A newly medium row has $d_i=0$. At its entry scale $r\le R$,
$u\le-H_0/r+a_0+b<-10$, and
\[
 2r^2w\le\frac{72}{r^2}e^{\beta(-H_0/r+a_0+b)}<B_0.
\]
The right side before the last inequality is increasing up to $R$, since
$\beta H_0/R>2$. At $R$ the inequality follows from
\begin{equation}\label{eq:initcertificate}
 \sum_{k=0}^{39}\frac{X^k}{k!}>\frac{72}{B_0R^2},
 \qquad X=\beta(H_0/R-a_0-b)=8.544644295\ldots.
\end{equation}
All quantities in this finite check are rational. Because
$pp^\T\preceq q\diag(p)$, a new medium row is covered in PSD order by
its former large-row deposit.

At fixed scale, deletion decreases $E$, $p$, $u$, and $w$, preserving $d$.
For decreasing scale at fixed row data,
\[
 r\partial_r u=b-u-E/r^2,\qquad
 r\partial_r\log w=\beta(b-u-E/r^2)-4.
\]
These are nonnegative as long as $u\le0$ and $r^2\ge q$, since
$\beta(b-a_0)=4$. The conditions persist during the decrease. Existing
medium Gram terms therefore decrease entrywise, which cannot increase the
top eigenvalue of a nonnegative symmetric matrix. Freezing removes a
principal submatrix; its removed energy terms vanish because $f(\pm1)=0$.
The vector $x$ itself is unchanged by cleanup. The column norm condition
and the deposit comparison give the initialization bound.
\end{proof}

\begin{lemma}[Tracking with the remaining mass retained]\label{lem:tracking}
Put $A_\Delta=(1+\sqrt{1-\Delta})/\sqrt\Delta
=(11+\sqrt{117})/2$.
At any time after a row enters the medium class, its deliberately deleted
absolute coefficient mass is at most
$A_\Delta(R-\sqrt{q_i})$.
For every later cube point $y$ agreeing with the frozen coordinates,
\begin{equation}\label{eq:tracking}
 (Ay)_i-\bigl(c_i+b_i^\T y_V\bigr)
 =\sum_{j\in D_i}a_{ij}(y_j-x_j^{\rm del}).
\end{equation}
\end{lemma}
\begin{proof}
At a deliberate deletion, $a^2>\Delta q$ and
\[
 \frac{|a|}{\sqrt q-\sqrt{q-a^2}}
 =\frac{\sqrt q+\sqrt{q-a^2}}{|a|}<A_\Delta.
\]
Sum these decreases of $\sqrt q$. Freezing only decreases $q$ further,
so the sum is at most $R-\sqrt{q_i}$. Offsets preserve the current tracked
sum at each removal. Telescoping the affine changes gives
\eqref{eq:tracking}; frozen removals make zero later contribution.
\end{proof}

\section{The coupled estimate and the dimension budget}\label{sec:coupling}

Put $z_i=p_i/r_i^2$ and $a_{i\sigma}=A_0e^{\beta u_{i\sigma}}$, so
$W=\sum a_{i\sigma}z_iz_i^\T$ and $\sum_jz_{ij}\le1$.
During differentiation retain the current row data and scales. Then
\[
 g_{i\sigma}=\nabla u_{i\sigma}
 =\sigma b_i/r_i+z_i\circ f'(x),\qquad
 \nabla^2u_{i\sigma}=\diag(z_i\circ f''(x)),\qquad
 \|g_{i\sigma}\|\le25/8.
\]
The last inequality follows from $\|z_i\|\le d$ and
$1+d\|f'\|_\infty=25/8$.

The ridge makes the Perron eigenvalue simple with gap at least $\eta$:
if $M_0=W+B$, a nonnegative unit top eigenvector $v_0$ has
$\one^\T v_0\ge1$, so $L\ge\lambda_1(M_0)+\eta$; min--max on
$\one^\perp$ gives $\lambda_2(M)\le\lambda_1(M_0)$.
Let $v>0$ be the unit Perron vector. Define
\begin{align}\label{eq:UD}
 \calU&=\sum_{i,\sigma}a_{i\sigma}\beta^2(z_i^\T v)^2
                                  g_{i\sigma}g_{i\sigma}^\T,\notag\\
 \calD&=\sum_{i,\sigma}a_{i\sigma}\beta(z_i^\T v)^2
                    \diag(z_i\circ(-f''(x))).
\end{align}
By \eqref{eq:pointwise},
\begin{equation}\label{eq:rho}
 \calU_{jj}\le\rho\calD_{jj},\qquad \rho=\frac{40}{187}.
\end{equation}

\begin{lemma}[Coupled low-spectrum response]\label{lem:coupled}
If $CC^\T\preceq M$, and $R_\alpha$ is the reduced resolvent of $M$
restricted to eigenvalues less than $\alpha L$, then
\[
 C^\T R_\alpha C\preceq\frac{\alpha}{1-\alpha}I.
\]
Consequently, writing $e_h=\|P_{[\alpha L,L]}M'[h]v\|$, the actual walk
satisfies
\begin{equation}\label{eq:curvature}
 L''[h,h]\le h^\T(\gamma\calU-\calD)h+2e_h^2/\eta,
 \qquad \gamma=\frac{1+\alpha}{1-\alpha}=\frac{309}{191}.
\end{equation}
\end{lemma}
\begin{proof}
The nonzero eigenvalues of $C^\T R_\alpha C$ equal those of
$R_\alpha^{1/2}CC^\T R_\alpha^{1/2}$, which is bounded above by
$R_\alpha^{1/2}MR_\alpha^{1/2}$. The latter has eigenvalues
$\lambda/(L-\lambda)<\alpha/(1-\alpha)$ on its support.
For the application let $C$ have columns $\sqrt{a_{i\sigma}}z_i$.
Then $M'[h]v=C\ell$, where
$\ell_{i\sigma}=\sqrt{a_{i\sigma}}\beta(z_i^\T v)g_{i\sigma}^\T h$,
and $\|\ell\|^2=h^\T\calU h$.
Also $v^\T M''[h,h]v=h^\T(\calU-\calD)h$.
Use the simple-eigenvalue second derivative formula, the displayed bound
on the low subspace, and the gap $\eta$ on the high subspace.
\end{proof}

For $S\succeq0$ and nonnegative diagonal $D$, the condition
$S_{jj}\le rD_{jj}$ implies that $S-D$ has fewer than $rs$ positive
eigenvalues, or none. Indeed, zero diagonal entries of $D$ annihilate the
corresponding rows of $S$; on the remaining support
$\tr(D^{-1/2}SD^{-1/2})\le rs$. Each eigenvalue above one uses more than
one unit of trace. This inertia estimate also holds after restriction to
any subspace.

There are fewer than $s/R^2$ large rows. Compute each barrier with absolute
error at most $1/1000$, and protect it if the computed value is at least
$-1/200$, imposing the exact equation $g_{i\sigma}^\T h=0$.
A protected pair has true $u\ge-1/100$, and therefore contributes at least
$A_0\theta^{-4}e^{-\beta/100}$ to $\one^\T W\one\le Ls$.
The number of protected pairs is at most
\begin{equation}\label{eq:protectedcount}
 \frac{L_*\theta^4}{36(1-\beta/100)}s.
\end{equation}
We used $e^t\le(1-t)^{-1}$ for $0\le t<1$.
Diffuseness also gives
\begin{equation}\label{eq:trace}
 \tr M\le(\Delta L+B_0+\eta)s,
\end{equation}
because $\|p_i\|^2\le\Delta q_i^2$ and
$\sum w_{i\sigma}q_i^2=\one^\T W\one\le Ls$.

Here is the full numerical slack, expressed through rational quantities.
Let $\ell_0=999999/1000000$. For $s\ge513$,
\begin{align}\label{eq:dimensioncertificate}
 1-\frac1{R^2}
 -\frac{L_*\theta^4}{36(1-\beta/100)}
 -\frac{\Delta+B_0/\ell_0}{\kappa}
 -\frac1{10000}
 -\frac{12360}{35717}-\frac1{513}
 &>\frac1{200}.
\end{align}
The left side is $0.00533725675760\ldots$.
The terms account for large rows, protected pairs, spectral cancellation,
numerical slack, positive curvature directions, and $x_V^\T h=0$,
respectively. The verifier checks this inequality with rational arithmetic.

\section{Computing an exactly feasible negative subspace}\label{sec:finite}

At phase start, with $S$ active coordinates, set
\begin{equation}\label{eq:precision}
 \eta=\frac1{2^{40}JS},\qquad \delta=\frac{\eta^5}{2^{100}}.
\end{equation}
Keep these fixed until at most $S/2$ coordinates remain. The active range
for this section is $s\ge513$.

Use the following finite symmetric-diagonalization interface: for a
bounded-norm symmetric $H$ and tolerance $\xi$, return a diagonal $D$ and
a matrix $V$ with an orthogonal analysis witness $U$ such that
\begin{equation}\label{eq:interface}
 \|V-U\|\le\xi,\qquad \|H-UDU^\T\|\le\xi.
\end{equation}
Rational Jacobi rotations implement it with $V=U$ in
$\widetilde O(s^3)$ field operations; the fast interface from the source
draft costs $\widetilde O(s^\omega)$ and permits square roots. Matrix norms
up to 16 are handled by scaling. Details of the arithmetic requirements
are collected in Section~\ref{sec:arithmetic}.

Approximate the nonnegative normalized weights by
$\bar a_{i\sigma}\ge0$ with
$\sum|\bar a_{i\sigma}-a_{i\sigma}|\le\delta/2^{20}$.
Form $\bar M$, its analytical derivative map $\bar M'[h]$, and apply
\eqref{eq:interface} with tolerance $\delta/2^{20}$. Clip reported negative
eigenvalues to zero. The resulting witness obeys
$\|M-UDU^\T\|<\delta$ and $\tr D\le\tr M+s\delta$.
Let $d_* =\max D_{jj}$ and let $w$ be the unit normalization of its
corresponding column of $V$, with fixed tie breaking. In rational Jacobi
the column already has unit length.

If $d_*\le1-\delta$, then $L\le1$: use the common kernel of the large-row
constraints, protected gradients, and $x_V^\T h=0$.
Otherwise $L>1-2\delta>\ell_0$. Let $V_F$ contain the columns whose
reported eigenvalues are at least $\kappa d_*$. Add the exact equations
\begin{equation}\label{eq:cancellation}
 V_F^\T\bar M'[h]w=0.
\end{equation}
Let $\mathcal H$ be the resulting exact kernel and $P$ its exact orthogonal
projector. The number of cancellation equations divided by $s$ is at most
\[
 \frac{\Delta+(B_0+\eta+\delta)/\ell_0}
      {\kappa(1-\delta/\ell_0)}
 <\frac{\Delta+B_0/\ell_0}{\kappa}+\frac1{10000}.
\]
The last inequality holds even with $\eta$ replaced by $2^{-40}$ and
$\delta$ replaced by $2^{-300}$, and is checked rationally.

\begin{lemma}[Numerical leakage and curvature]\label{lem:leakage}
In the high branch, orient $w$ for analysis so that
$\|w-v\|\le3\delta/\eta$, with $v>0$.
This orientation changes neither the equations nor the proxy.
Every $h\in\mathcal H$ satisfies
\begin{equation}\label{eq:leakage}
 \|P_{[\alpha L,L]}M'[h]v\|,\ |L'[h]|
 \le\frac{2^{20}\delta}{\eta}\|h\|.
\end{equation}
Define $\bar{\calU},\bar{\calD}$ by \eqref{eq:UD}, replacing weights by
$\bar a$ and $v$ by $w$, and put $\bar C=\gamma\bar{\calU}-\bar{\calD}$.
Then $\|\bar C\|<16$, its positive inertia is less than
$(12360/35717)s$, and
\[
 \|\gamma\calU-\calD-\bar C\|\le2^{12}\delta/\eta.
\]
Every feasible $h$ with $h^\T\bar C h<\eta\|h\|^2$ obeys
\begin{equation}\label{eq:potentialcurvature}
 \Phi''[h,h]\le-\frac{19}{10}\eta\|h\|^2,
 \qquad |\Phi'[h]|\le\frac{2^{20}\delta}{\eta}\|h\|.
\end{equation}
\end{lemma}
\begin{proof}
The gap $\eta$ and the residual of the computed top column give the stated
eigenvector estimate. Between true eigenvalues at least $\alpha L$ and
witness eigenvalues below $\kappa d_*$, separation is at least
$(\alpha-\kappa)L-\kappa\delta>1/2000$.
The Sylvester equation for these two spectral blocks therefore gives
$\|P_{[\alpha L,L]}U_{F^c}\|\le2000\delta$.
Also $\|M'[h]\|\le(5/2)L\|h\|<3\|h\|$ and
$\|\bar M'[h]-M'[h]\|\le(5/2)\delta\|h\|/2^{20}$.
Replace $v$ by $w$, use \eqref{eq:cancellation}, and replace $V_F$ by
$U_F$. These estimates yield \eqref{eq:leakage} with room to spare.

The diagonal domination \eqref{eq:rho} holds for any vector $w$ and
nonnegative weights, so the inertia statement is exact.
The bounds $\beta\|g\|\le5/2$,
$\beta\|f''\|_\infty\Delta<1$, and
$\sum\bar a(z_i^\T w)^2\le L+\delta$ give $\|\bar C\|<16$.
For the perturbation use
\[
 \sum a\left|(z_i^\T v)^2-(z_i^\T w)^2\right|
 \le2L\|v-w\|.
\]
Each coefficient matrix in $\gamma\calU-\calD$ has norm less than 12;
including the weight error proves the stated bound.

Combine it with \eqref{eq:curvature} and \eqref{eq:leakage}. For
$F(t)=(t-1)_+^3$, $F'(L)\le3/65536$ and $F''(L)\le3/128$ at admissible
states. The second derivative is
$F'(L)L''+F''(L)(L')^2-2\eta\|h\|^2$.
Using \eqref{eq:precision} leaves the bound in
\eqref{eq:potentialcurvature}. The first derivative of the quadratic
energy is zero because $x_V^\T h=0$. When $L\le1$, the spectral first and
second derivatives vanish.
\end{proof}

In the high branch form
$Q_{\rm neg}=P(\bar C-\eta I)P+(I-P)$; in the low branch form
$Q_{\rm neg}=-\eta P+(I-P)$.
By \eqref{eq:dimensioncertificate}, $Q_{\rm neg}$ has more than $s/200$
eigenvalues at most $-\eta$, and $\|Q_{\rm neg}\|<16$.
Apply \eqref{eq:interface} with error $\xi=\eta/2^{20}$ and retain
columns reported at most $-3\eta/4$, giving $V_-$ with $k>s/200$ columns.
Set
\begin{equation}\label{eq:basis}
 Z=PV_-,\qquad T=Z/(1+\xi).
\end{equation}
Then $T$ is exactly feasible and
$\frac34 I_k\preceq T^\T T\preceq I_k$.
Indeed the orthogonal witness satisfies
$(I-P)U_-(I-D_-)=(I-P)(Q_{\rm neg}-UDU^\T)U_-$,
so $\|(I-P)U_-\|\le\xi$ and $\|Z-U_-\|\le2\xi$.
The computed span is negative for $Q_{\rm neg}$: the error from replacing $U_-$ by
$V_-$ is at most $34\xi<\eta/4$.
Projection preserves negativity because
\[
 (Za)^\T(\bar C-\eta I)(Za)
 =(V_-a)^\T Q_{\rm neg}(V_-a)-\|(I-P)V_-a\|^2<0.
\]
In the low branch all feasible vectors already have spectral curvature zero.

\section{Flat directions and finite steps}\label{sec:steps}

Apply the deterministic flattening construction to the test vectors
$e_j$ and $(8/25)g_{i\sigma}$, each of norm at most one, and to the matrix
$T$ from \eqref{eq:basis}. It returns a rational vector $h_0$ in its span
with
\[
 \frac34\le\|h_0\|\le1,\qquad
 |(h_0)_j|\le3\sqrt{q/k},\qquad
 |g_{i\sigma}^\T h_0|\le\frac{75}{8}\sqrt{q/k}.
\]
Put $h=h_0/16$, so $1/32\le\|h\|\le1/16$.
If $4q/k\ge1$ use $K=1$; otherwise halve from one until the next square
would fall below $4q/k$. Then
\begin{equation}\label{eq:flat}
 |\beta g_{i\sigma}^\T h|\le K,\quad
 8\|h\|_\infty\le K,\quad
 \eta\le K\le1,\quad K^2\le\min(1,3200q/s).
\end{equation}
The baseline inequalities when $K=1$ follow directly from $\|h\|\le1/16$.
In the other case the displayed flat estimates suffice.
Section~\ref{sec:arithmetic} records the finite conditional-expectation
implementation, so no random or continuous oracle is used here.

Along a cube-contained line $x+th$, write
$w_{i\sigma}(t)=w_{i\sigma}(0)e^{\ell_{i\sigma}(t)}$.
The derivative bounds \eqref{eq:fderivatives} and \eqref{eq:flat} imply
\begin{equation}\label{eq:logcurve}
 \ell(0)=0,\quad |\ell'(0)|\le K,\quad
 -K^2/2\le\ell''(t)\le0,\quad
 |\ell'''(t)|\le K^3/4\quad\text{a.e.}
\end{equation}
The last two assertions use $\sum z_{ij}\le1$ and
$8\|h\|_\infty\le K$. They remain valid when coordinates cross zero.

\begin{lemma}[Finite Perron comparison for $C^{2,1}$ curves]\label{lem:perron}
Let $M(t)=B+\sum_r w_re^{\ell_r(t)}p_rp_r^\T$, with $B\succeq0$
entrywise nonnegative, $p_r\ge0$, and $w_r\ge0$.
Assume \eqref{eq:logcurve}, $K\le1$, and at zero assume
$1/2\le L\le9/8$, gap at least $\zeta\in(0,1]$, and unit Perron vector
$v>0$. Write $\nu=\min v_i$, $k_v=1+\log(1/\nu)$,
$E=M'(0)$, and suppose
\[
 e:=\|P_{(L/4,L]}Ev\|\le K\zeta\nu/64.
\]
For $|t|\le\sqrt\zeta/(64Kk_v)$ within the available curve interval,
\begin{equation}\label{eq:remainder}
 \Phi(t)\le\Phi(0)+\Phi'(0)t+\tfrac12\Phi''(0)t^2
             +128K^3|t|^3+4096K^4t^4/\zeta,
\end{equation}
where $\Phi=(\lambda_1(M)-1)_+^3$ minus any quadratic energy.
\end{lemma}
\begin{proof}
The step restriction gives $K|t|\le1/64$. From \eqref{eq:logcurve} and the
positive Gram representation,
\[
 |E_{ij}|\le KM(0)_{ij},\quad \|M''(t)\|\le6K^2,\quad
 \|M'''(t)\|\le14K^3\quad\text{a.e.}
\]
Let $\alpha_0=v^\T Ev$, let $R_0$ be the reduced resolvent, and put
$z=R_0(Ev-\alpha_0v)$.
Splitting the forcing at $L/4$ gives $\|z\|\le2K$ and
$|\alpha_0|\le e\le K/64$.
For the low response use the series
$L^{-1}\sum_{j\ge0}(M(0)/L)^j$ on the low space.
Its $j$th component relative to $v_i$ is bounded both by $(33/32)K$
(positivity applied to $Ev$, plus the small removed high part) and by
$K4^{-j}/\nu$ (spectral decay). Summing the smaller bounds, and adding
the high response $e/(\zeta\nu)$, gives
\begin{equation}\label{eq:zbound}
 |z_i|\le4Kk_vv_i.
\end{equation}

Put $A=LI-M(0)$, $Z=\diag(z_i/v_i)$, and $D=I+tZ$.
Then $Az=Ev-\alpha_0v$ and
\[
 D((L+\alpha_0t)I-M(t))D=A+tG+R(t),\quad
 G=ZA+AZ-E+\alpha_0I,
\]
where $Gv=0$ and $\|R(t)\|\le80K^2k_v^2t^2$.
This follows by expanding, using $\|A\|\le2$,
$\|E\|\le2K$, \eqref{eq:zbound}, and
$\|M(t)-M(0)-tE\|\le6K^2t^2$.
Off the diagonal $|G_{ij}|\le9Kk_vM(0)_{ij}$.
For a symmetric matrix annihilating $v$, the identity
\[
 y^\T Gy=-\frac12\sum_{ij}G_{ij}v_iv_j
                         (y_i/v_i-y_j/v_j)^2
\]
therefore implies $|y^\T Gy|\le9Kk_v y^\T Ay$.
Since $\|D-I\|\le1/16$, min--max on $(Dv)^\perp$ gives
\[
 \lambda_2(M(t))\le L+\alpha_0t-\zeta/2.
\]
Explicitly, the lower bound for the complementary quadratic form is
\[
 \left[\frac{55}{64}\left(\frac{16}{17}\right)^2
 -\frac{80}{4096}\left(\frac{16}{15}\right)^2\right]\zeta
 >\zeta/2.
\]

The positive trial vector $\psi=v+tz$ has Rayleigh quotient $r(t)$.
Its residual centered at $L+\alpha_0t$ is
\[
 (M(t)-(L+\alpha_0t)I)\psi
 =t^2(E-\alpha_0I)z+(M(t)-M(0)-tE)\psi.
\]
Its normalized norm is at most $13K^2t^2$.
Thus $|r-L-\alpha_0t|\le13K^2t^2$ and
$r-\lambda_2(M(t))>\zeta/3$.
The elementary Rayleigh residual (Temple) inequality gives
\begin{equation}\label{eq:temple}
 0\le\lambda_1(M(t))-r(t)\le512K^4t^4/\zeta.
\end{equation}

Only a second-order Taylor expansion is needed for the $C^{2,1}$ curve.
Put $N=M''(0)$; its remainder has norm at most
$(14/6)K^3|t|^3$. Expand
$(v+tz)^\T M(t)(v+tz)/(1+t^2\|z\|^2)$.
The quadratic coefficient is $L''(0)/2$ because
$z^\T M(0)z=L\|z\|^2-z^\T Ev$.
The remaining cubic coefficient is bounded using
$\|z\|\le2K$, $\|E\|\le2K$, $\|N\|\le6K^2$,
and $|\alpha_0|\le K/64$; the quartic terms are absorbed using
$K|t|\le1/64$. This gives explicitly
\[
 \left|r(t)-L-\alpha_0t-\tfrac12L''(0)t^2\right|
 \le32K^3|t|^3.
\]
For completeness, the numerator's cubic terms have bound
$(8+12+1/16)K^3$, its quartic terms have bound $40K^4$, and the
normalized Taylor remainder has bound $(14/6)K^3|t|^3$.

For $F(u)=(u-1)_+^3$, $F''$ is globally 6-Lipschitz, so its scalar
second-order remainder is at most $|r-L|^3$.
Use $|r-L-\alpha_0t|\le13K^2t^2$ and
$|\alpha_0|\le K/64$ to get a bound of $128K^3|t|^3$ after composition.
All arguments lie below two. Replacing $r$ by $\lambda_1$ using
\eqref{eq:temple} costs at most $1536K^4t^4/\zeta$, which is smaller
than the final allowance in \eqref{eq:remainder}. A quadratic energy
contributes no remainder. The proof uses no fourth derivative.
\end{proof}

In the high branch, $Lv_i\ge\eta\sum_jv_j\ge\eta$, so
$\nu\ge\eta/2$. Since $\alpha<1/4$, \eqref{eq:leakage} controls the
high space in Lemma~\ref{lem:perron}.
Equations \eqref{eq:precision} and \eqref{eq:flat} imply its leakage
condition with $\zeta=\eta$.
In the low branch $L\le1$, and the spectral potential after a short move
is bounded directly by $64K^3|t|^3$; its first two derivatives at zero vanish.

Set $c_*=2^{-40}$ and let $H=\lceil\log_2(4/\eta)\rceil$.
Take the largest dyadic $\tau\le1$ satisfying
\begin{equation}\label{eq:stepschedule}
 \tau\le c_*,\qquad K^2\tau\le c_*\eta,\qquad
 K^2H^2\tau^2\le c_*^2\eta.
\end{equation}
These are rational tests. We have $H\ge k_v$, $\eta H<1$,
and $\tau\ge c_*\eta/(2H)$.
The precision choice implies
$|\Phi'[h]|\le\eta\tau/2^{16}$, and
\eqref{eq:potentialcurvature} implies $\Phi''[h,h]\le-\eta/1024$.
Lemma~\ref{lem:perron} applies on $|t|\le\tau$.
Its remainder, divided by $\eta t^2$, is at most
$128c_*+4096c_*^2<2^{-16}$.
For either sign and a cube-contained move $0<t\le\tau$,
\begin{equation}\label{eq:onestep}
 \Phi(x\pm th)-\Phi(x)
 \le\frac{\eta\tau t}{2^{16}}-\frac{\eta t^2}{2048}
                              +\frac{\eta t^2}{2^{16}}.
\end{equation}
A full step decreases $\Phi$ by at least $\eta\tau^2/4096$.
A shorter step increases it by at most $\eta\tau^2/2^{16}$.
Protected barriers cannot increase because they are concave and tangent to
$h$. Every unprotected barrier is below $-1/250$ and its increase is at
most $\beta^{-1}K\tau<1/250$. Large rows remain zero.

\section{The complete algorithm and its spectral budget}\label{sec:algorithm}

All tie choices below are by least index, with the positive sign preferred.
First perform the row preprocessing of Section~\ref{sec:screening}.
For $n=0$ return the empty signing; if $m_0=0$, return all positive signs.
For $1\le n\le512$ go directly to the
terminal procedure of Section~\ref{sec:terminal} at $x=0$.
Otherwise initialize the retained-row state and the first phase at $S=n$.
Repeat:
\begin{enumerate}
\item Compute the approximate weights and protected pairs, the exactly
feasible matrix $T$, the flat direction $h$, $K$, and $\tau$ as above.
\item Let $t_*=\min_{h_j\ne0}(1-|x_j|)/|h_j|$ and
$t=\min(t_*,\tau)$. If $t=\tau$, move by $+th$.
If $t<\tau$, choose a coordinate attaining $t_*$ and choose the move sign
so that this coordinate moves toward its nearer boundary (either sign if
it is zero). Move $x_V\leftarrow x_V\pm th$.
\item Perform cleanup. If $s\le512$, stop the walk. If $s\le S/2$, begin
a new phase with $S=s$ and the new parameters \eqref{eq:precision}.
\end{enumerate}
Then perform terminal conditional-expectation rounding.
Both signs of a move of size at most $t_*$ remain in the cube, and a short
move freezes a coordinate. The algorithm never evaluates the true
potential or an exact eigenvector.

\begin{lemma}\label{lem:budget}
All states obey $L<L_*$, and there are $O(nJ^2q^2)$ updates.
\end{lemma}
\begin{proof}
Write $F=(L-1)_+^3$. Within a phase, at most $S$ short steps occur and
the increase in $\|x\|^2$ is at most $S$. Equation~\eqref{eq:onestep}
therefore bounds every phase prefix's increase in $F$ by
$(1+2^{-16})\eta S=(1+2^{-16})/(2^{40}J)$.
A reset changes $L$ by at most $(\eta'-\eta)S'\le1/(2^{40}J)$.
If the preceding state has $L<L_*$, the tentative reset value is below
$1+1/128$, where $F'<1$. Thus a reset increases $F$ by at most
$1/(2^{40}J)$. There are at most $J$ phases and fewer than $J$ resets,
so every prefix has
\[
 F<3\cdot2^{-40}<2^{-24}=(L_*-1)^3.
\]
This inductively proves the invariant. Equivalently, at a first putative
crossing the local step bounds already give the same strict inequality,
so the argument does not assume the next state's admissibility.

During a phase $S/2<s\le S$. By \eqref{eq:flat},
$\eta/K^2>1/D_*$, where $D_*=2^{41}\cdot3200Jq$.
For $n\ge513$, $q\ge J$, $H\le6J$, and $D_*>H^2$.
Thus the largest dyadic step satisfies $\tau\ge c_*/(2D_*)$.
Every full displacement has norm at least $\tau/32=\Omega((Jq)^{-1})$.
The exact constraint $x_V^\T h=0$ gives
$\|x\pm th\|^2-\|x\|^2=t^2\|h\|^2$.
The total increase is at most $n$, so only $O(nJ^2q^2)$ full updates occur.
There are at most $n$ short updates. A phase beginning at size $S$ has
$O(SJ^2q^2)$ updates by the same argument.
\end{proof}

\section{Terminal rounding from the remaining allowance}\label{sec:terminal}

At $s\le512$, round active coordinates to signs with
\begin{equation}\label{eq:terminalenergy}
 \|A_V(\eps_V-x_V)\|_2^2+4\|\eps_V-x_V\|_2^2\le5s.
\end{equation}
This is deterministic polynomial-time conditional expectation.
Independent signs with means $x_j$ have expected left side
\[
 \sum_{j\in V}(1-x_j^2)(\|A_{:j}\|_2^2+4)\le5s.
\]
Fix one sign at a time, selecting the smaller of its two conditional
quadratic expectations; these are rational expressions in the current
data. Maintaining the partial vector $A_V(\eps_V-x_V)$ gives $O(ms+s)$
arithmetic work. In particular,
\begin{equation}\label{eq:roundnorms}
 \|A_V(\eps_V-x_V)\|_\infty\le\sqrt{2560},\qquad
 \|\eps_V-x_V\|_2\le\sqrt{640}.
\end{equation}

A still-large row has actual current sum zero and has had no deliberate
deletions. Its final discrepancy is at most $\sqrt{2560}<C_*$.
For a medium row, its two barriers imply
$|d_i|\le H_0-br_i\le H_0-b\sqrt{q_i}$.
Lemma~\ref{lem:tracking} and Cauchy--Schwarz on its retained coefficients
give
\begin{align*}
 |(A\eps)_i|
 &\le H_0-b\sqrt{q_i}
        +\sqrt{640}\sqrt{q_i}
        +2A_\Delta(R-\sqrt{q_i})\\
 &=H_0+2RA_\Delta
       +\bigl(\sqrt{640}-b-2A_\Delta\bigr)\sqrt{q_i}
 <H_0+2RA_\Delta.
\end{align*}
The last coefficient is negative because
$b+2A_\Delta=30.536653826\ldots>\sqrt{640}$.
A completed row has $|d_i|<H_0$ (or $d_i=0$ if it went directly from large
to completed), and its deletion error is bounded by $2RA_\Delta$.
If $n\le512$, apply \eqref{eq:terminalenergy} at zero; its discrepancy is
at most $\sqrt{5n}\le\sqrt{2560}<C_*$.
Finally,
\[
 H_0+2RA_\Delta=H_0+R(11+\sqrt{117})
 =\frac{1271}{25}+\frac{149}{100}\sqrt{117}<67.
\]
The last inequality is certified without square-root approximation by
$(67-1271/25)^2>117(149/100)^2$.
This proves the discrepancy assertion of Theorem~\ref{thm:main}.

\section{Finite arithmetic and implementation scope}\label{sec:arithmetic}

\subsection{Profile and scalar evaluations}
The row constraints use $f'$ and the proxy uses $f''$, both rational in
rational $x$. Only weights and protected-pair decisions require values of
$f$. For $z\in[1,25/8]$, let $t=(z-1)/(z+1)\in[0,17/33]$. Then
\[
 \log z=2\sum_{k=0}^{N-1}\frac{t^{2k+1}}{2k+1}+R_N,
 \qquad 0\le R_N\le\frac{2t^{2N+1}}{(2N+1)(1-t^2)}.
\]
Thus $O(\log(1/\xi))$ rational operations give any absolute accuracy
$\xi$. Compute the profile at each active coordinate, then form the row
energies. Since $\sum_jz_{ij}\le1$, profile error is not magnified in
$E_i/r_i^2$. The other term $(\sigma d_i-H_0)/r_i$ is exact arithmetic.
Compute positive weight approximants with the total error required in
Section~\ref{sec:finite}. Very negative exponentials can be set to zero
after certifying that their weight is below the allocated error.
For the rest, argument reduction and a positive Taylor series suffice.
The needed precision logarithms are $O(\log(2+m+n))$. No logarithm or
exponential oracle is assumed. Tiny scales cause no loss of the absolute
weight accuracy because the matrix is formed with $a_{i\sigma}z_iz_i^\T$,
not with unnormalized $r_i^{-4}$ coefficients.

\subsection{An explicit field-operation diagonalization interface}
For a symmetric matrix of bounded norm maintain an exactly orthogonal
product of rational plane rotations. In a pivot block with diagonal
$a,c$ and off-diagonal $b$, use
$\cos t=(1-z^2)/(1+z^2)$ and $\sin t=2z/(1+z^2)$.
The rotated off-diagonal numerator is
$b(1-6z^2+z^4)+2z(1-z^2)(c-a)$.
A zero lies between $0$ and
$-\tfrac12\operatorname{sgn}(b(c-a))$; if $a=c$, use the endpoint $1/2$.
Bisection reduces the pivot magnitude by at least a factor two.
Choosing a largest pivot reduces squared off-diagonal Frobenius norm
$E$ by at least $3E/(2s(s-1))$.
After $O(s^2\log(s/\xi))$ rotations the off-diagonal norm is below $\xi$.
A tournament structure updates the largest pivot in $O(s\log s)$ work
per rotation. While termination is not yet reached, the pivot is at least
an inverse-polynomial multiple of $\xi$, so its bisection takes
$O(\log(s/\xi))$ iterations. This proves
\eqref{eq:interface} in $\widetilde O(s^3)$ arithmetic operations.
It also handles indefinite $Q_{\rm neg}$.

Extract independent exact constraint rows into $E_0$ and form
$P=I-E_0^\T(E_0E_0^\T)^{-1}E_0$; use $P=I$ if there are no rows.
Exact elimination and projection cost $O(s^3)$. This model allows
arbitrarily small nonzero pivots and makes no conditioning assumption.

\subsection{Finite deterministic flattening}
Here are details of the flattening step used above, to specify the
algorithm without a probabilistic oracle. For unit-norm test vectors
$a_1,\ldots,a_N$, write $c_{ij}=a_i^\T T e_j$; then
$\sum_jc_{ij}^2\le1$.
Let $\lambda$ be the least positive integer with $\lambda^2\ge q$.
For uniform independent signs $\eps_j$,
\[
 \mathbb E\sum_i\cosh\!\left(\lambda\sum_jc_{ij}\eps_j\right)
 \le N e^{\lambda^2/2}.
\]
Greedily fix each sign using positive relative approximations to the two
conditional expectations. Taking scalar relative errors at most
$1/(100k^2)$ makes each conditional expectation accurate within
$1/(20k)$ and the accumulated multiplicative loss less than two.
The resulting signs have
$|\sum_jc_{ij}\eps_j|\le3\sqrt q$ for every $i$:
use $\log(4N)+\lambda^2/2\le3q$ and $\lambda\ge\sqrt q$.
Choose a positive dyadic rational $r_0$ with
$3/4\le r_0^2k\le1$, by bisection, and put $h_0=r_0T\eps$.
The Gram bounds on $T$ give $3/4\le\|h_0\|\le1$ and the claimed test bounds.
The two conditional expectations are positive sums of products of
$\cosh(\lambda c_{ij})$ and a current partial-sum factor.
Suffix products and partial sums yield $O(Nk)$ scalar evaluations.
Arguments have size at most $2\sqrt{qk}$; halving, Taylor approximation,
and repeated squaring implement their relative approximants in
$O(\log(2Nk))$ operations each.

\subsection{Unscreened core cost and implementation scope}
The per-state field-operation cost is
$\widetilde O(ms^2+s^3+m)$, including Gram and proxy formation, both
diagonalizations, exact kernel computation, and flattening.
Each phase beginning at $S$ has $O(SJ^2q^2)$ updates.
Summing the geometric phase sizes gives
$\widetilde O(m_0n^3+n^4)$ after the $O(mn)$ input scan.
The extra profile evaluations and terminal rounding fit within this bound.

The source draft's fast diagonalization interface satisfies
\eqref{eq:interface} with an orthogonal analysis witness, while fast exact
rank computation constructs the same projector. Blocking rectangular
products by $s$ rows makes Gram, proxy, cancellation, and flattening
products cost $O(ms^{\omega-1}+s^\omega)$ per state.
The same phase sum yields
$\widetilde O(m_0n^\omega+n^{\omega+1})$ in that arithmetic model.
This is the cost of evaluating every surviving row on every update.
Section~\ref{sec:screening} removes that repeated cost and proves
\eqref{eq:screenedmaincost}. Neither version analyzes the separate
randomized and amortized refinements in the longer source draft.

The step constants are deliberately conservative. This is a complete
finite arithmetic algorithm and proof, with an executable exact parameter
verifier, rather than a benchmarked end-to-end numerical solver. Exact
rank computation may increase bit lengths; no polynomial bit bound is
inferred from the arithmetic count. The separate existence bound near
$6.9013$ remains substantially smaller than the constructive guarantee.

\section{Certificates and reproducibility}

The accompanying \texttt{verify\_constant67.py} uses only Python's standard
library. It certifies the profile height by a rational logarithm enclosure,
the initialization inequality by a finite positive Taylor sum,
the entire dimension budget including numerical slack, the terminal
rounding inequalities, and the finite-step coefficient bounds.
It also verifies $C_*<67$ using a rational square comparison. Its decimal
output is for readability; the assertions use exact fractions.
The proof of the $C^{2,1}$ finite-step lemma is given above, including the
zero-crossing case, and does not rest on a floating-point experiment.
The optional \texttt{audit\_finite\_profile.py} additionally checks that
lemma on noncommuting Gram curves crossing the third-derivative kink,
with two nearly colliding leading eigenvalues and gaps down to order
$10^{-24}$, using 140-digit arithmetic. Those checks are supporting audits,
not substitutes for the proof or end-to-end solver tests.
The additional standard-library script
\texttt{verify\_screened\_complexity.py} (not included in this repository)
certifies the screened dimension
margin, checks the two aggregate statistics through signed
coefficient deletion and freezing using exact fractions, and checks cold-row
clock certificates through movement and support changes, and verifies
the exact row-energy increment along doubly tangent directions. It also checks
the screening precision schedule on several input sizes. These tests
support the symbolic analysis below; they do not claim empirical runtime
exponents.

\section{Screening rows and the improved complexity}\label{sec:screening}

This section proves the running time in Theorem~\ref{thm:main}.
The signing bound and its constants are unchanged. The modifications are
an initial row filter, direct row-energy control, two extra sets of direction
constraints, and certified omission of very small Gram weights.

\subsection{One scan removes the dependence on the original row count}
A row with $\|a_i\|_1\le66$ is safe for every signing, because
$|(A\eps)_i|\le66<C_*$. Retain just the other $m_0$ rows. For every
survivor, Cauchy--Schwarz gives $\|a_i\|_2^2>66^2/n$; summing the column
norm assumptions gives $\sum_i\|a_i\|_2^2\le n$. Thus
\begin{equation}\label{eq:rowpreprocess}
 m_0<n^2/66^2.
\end{equation}
The scan costs $O(mn)$ field operations and comparisons. The surviving
matrix still has column norms at most one. Terminal rounding and all
subsequent row operations use this matrix only; the omitted rows require
no subsequent verification.

This observation alone turns the unscreened fast bound into
$O(mn)+\widetilde O(n^{\omega+2})$. Achieving the stronger
\eqref{eq:screenedmaincost} requires the following changes.

\subsection{Controlling row energy with the original potential}
Let $\mathsf A$ be the matrix of current retained rows $b_i^\T$, and let
$c$ be the vector of offsets, including completed rows. Define
\[
 Q=\mathsf A^\T\mathsf A,\qquad u=\mathsf A^\T c,\qquad
 \mathcal E=\|c+\mathsf A x_V\|^2=\sum_i d_i^2.
\]
Then $\tr Q\le s$ and $\|Q\|\le s$.
Keep the original potential $\Phi=(L-1)_+^3-\eta\|x\|^2$ and phase
schedule. Write $\eta_0=1/(2^{40}Jn)$ for its minimum ridge parameter.
At each direction computation impose, in addition to all earlier exact
constraints,
\begin{equation}\label{eq:chargegradient}
 (Qx_V+u)^\T h=0.
\end{equation}
This makes the first derivative of $\mathcal E$ zero.

Also apply the symmetric-diagonalization interface to $Q$, requesting
original matrix and witness error at most $\xi_Q/4$, where
$\xi_Q=1/(2^{20}(s+1))$. Its norm is handled by scaling by $s+1$.
Clip negative reported eigenvalues to zero and impose $V_Q^\T h=0$ for
the computed columns whose reported eigenvalues are at least 4095.
Because $Q$ is PSD, every negative reported eigenvalue has magnitude at
most $\xi_Q/4$. Clipping therefore leaves total matrix error below
$\xi_Q$, while the witness error is also below $\xi_Q$. Thus the clipped trace
is at most $s+s\xi_Q$. The number of new columns is less than $s/4094$.
For every vector in their exact computed kernel,
\begin{equation}\label{eq:chargequadratic}
 h^\T Qh\le4096\|h\|^2.
\end{equation}
Indeed, its component along the corresponding orthogonal witness columns
is at most $\xi_Q\|h\|$. On the complementary witness subspace the
reported eigenvalues are below 4095; the true quadratic form is therefore
at most $[4095+(s+\xi_Q)\xi_Q^2+\xi_Q]\|h\|^2<4096\|h\|^2$.

The extra equation \eqref{eq:chargegradient} and these columns leave enough
room in \eqref{eq:dimensioncertificate}. Let $d_{\rm core}$ denote the
rational left side of that inequality. The exact certificate gives
\begin{equation}\label{eq:chargeddimension}
 d_{\rm core}-\frac1{513}-\frac1{4094}
 =0.00314367912628\ldots>\frac1{400},
\end{equation}
where the decimal is reported only for readability.
Construct the negative basis exactly as before, using this smaller exact
kernel. It now has $k>s/400$ columns. The same deterministic flattening
and normalization give
\begin{equation}\label{eq:screenedflat}
 \frac1{32}\le\|h\|\le\frac1{16},\qquad
 K^2\le\min(1,6400q/s).
\end{equation}
Only the test forms specified below need to be flattened.

No new curvature term is added to the potential, so the earlier
$-(19/10)\eta\|h\|^2$ margin is unchanged. The point of the new
constraints is the exact increment identity
\begin{align}\label{eq:rowenergyincrement}
 \mathcal E(x+th)-\mathcal E(x)
 &=t^2h^\T Qh\notag\\
 &\le4096t^2\|h\|^2
 =4096\bigl(\|x+th\|^2-\|x\|^2\bigr).
\end{align}
The first equality uses \eqref{eq:chargegradient}; the last uses the
original exact constraint $x_V^\T h=0$.
Both identities hold for either move sign. Initially $d=0$ and $x=0$,
and cleanup preserves every $d_i$ and the full vector $x$. Phase resets
change neither. Therefore, throughout the walk,
\begin{equation}\label{eq:rowenergybound}
 \sum_i d_i^2=\mathcal E\le4096\|x\|^2\le4096n.
\end{equation}
This is an algebraic consequence of the chosen directions. It requires
no extra term in the spectral potential and no additional progress
estimate. It is used only before terminal rounding.

\subsection{Which rows are evaluated}
Let $\delta_0=\eta_0^5/2^{100}$ and set
\begin{equation}\label{eq:screeningprecision}
 \rho_c=\frac{\delta_0}{2^{22}(m_0+1)},\qquad
 2^{\ell_c}\ge\frac{36}{\rho_c},\qquad
 C_c=a_0+b+\ell_c/\beta,
\end{equation}
where $\ell_c$ is the least nonnegative integer satisfying the displayed
test. These are rational constructions. By \eqref{eq:rowpreprocess},
$\ell_c,C_c=O(\log(2+n))$.
A medium row has score
\[
 \chi_i=|d_i|+C_c r_i.
\]
If $\chi_i\le H_0$, then
$u_{i\sigma}\le-\ell_c/\beta$ for both signs, so its normalized weights
$a_{i\sigma}=36e^{\beta u_{i\sigma}}$ are at most $\rho_c$.
It has no protected pair. A row is called cold if a certificate proves
$\chi_i\le H_0$ at the current anchor; otherwise it is warm.

Use hysteresis as follows. Refresh all current warm rows before a move;
demote one only when $\chi_i\le H_0-2$. On checking a cold row, renew its
certificate if $\chi_i\le H_0-1$, and otherwise promote it. Check a newly
medium row once, with the latter rule. Expired certificates are processed
before computing the direction. Section~\ref{sec:clock} gives the exact
certificate and its expiration time.

In numerical matrix formation, set cold weights to zero. Approximate all
warm weights with total absolute error at most $\delta/2^{22}$.
The sum of omitted cold weights is less than
$2m_0\rho_c<\delta_0/2^{21}\le\delta/2^{21}$.
Thus the total error remains below the original allowance
$\delta/2^{20}$. The computed spectral kernel and curvature proxy are
formed from warm weights. The flat test list consists of the coordinate
vectors and the warm medium gradients. Large-row constraints, the exact
row-energy constraint, and the $Q$ constraints are retained.

\subsection{Cold rows do not spoil the larger finite steps}\label{sec:coldstep}
It is essential to justify removing cold gradients from the flat test
list: their unweighted directional derivatives need not satisfy the small
$K$ bound. We use a comparison that only requires their tiny total weight.
At a fixed anchor write the full true Gram contribution as
$W=W_{\rm warm}+W_{\rm cold}$. For the analysis of this one move, define
\[
 M_{\rm fr}(t)=B+\eta\one\one^\T+W_{\rm cold}(0)+W_{\rm warm}(t).
\]
This is an analysis curve, not a matrix computed by the algorithm.
At zero $M_{\rm fr}(0)=M(0)$, so its Perron eigenvector and gap are those
of the full matrix. The fixed cold contribution is PSD and entrywise
nonnegative. Every varying weight is warm and satisfies
\eqref{eq:logcurve}; hence Lemma~\ref{lem:perron} applies to this curve.
The coupled-response proof also applies: its Gram factor uses just the
warm terms and is still dominated by the full anchor matrix.
The proxy's diagonal comparison, spectral trace count, and numerical
leakage bounds are unchanged. Omitting the cold derivatives introduces
no derivative error relative to this particular frozen-cold curve.

For a cold pair on a cube-contained move, the global gradient bound gives
$|\beta g^\T h|\le5/32$. Also
$|\ell''(t)|\le K^2/2\le1/2$ from the coordinate flatness, which holds
for all rows. Consequently, for $|t|\le c_*=2^{-40}$,
$|\ell(t)|\le(5/32)|t|+t^2/4<|t|/4$, and
$|e^{\ell(t)}-1|\le|t|$. Summing the positive Gram factors gives
\begin{equation}\label{eq:coldmotion}
 \|M(t)-M_{\rm fr}(t)\|\le\frac{\delta_0}{2^{21}}|t|.
\end{equation}
The two spectral arguments lie below two on these steps, and
$z\mapsto(z-1)_+^3$ is 3-Lipschitz there. Their spectral potentials
therefore differ by at most $3\delta_0|t|/2^{21}$.
The quadratic term $-\eta\|x+th\|^2$ is identical on both curves.

Apply the frozen-cold version of \eqref{eq:potentialcurvature} and use
the exact tangency $x_V^\T h=0$.
The former first-order error and the additional term from
\eqref{eq:coldmotion} are together at most
$(2^{21}\delta/\eta)|t|$.
The unchanged precision and step schedules give
$2^{21}\delta/\eta\le\eta\tau/2^{16}$:
use $\tau\ge2^{-41}\eta/H$, $\eta H<1$, and
$\delta=\eta^5/2^{100}$.
It follows that \eqref{eq:onestep} still holds for the original $\Phi$.
Protected and warm unprotected barriers remain feasible by the old
argument. A cold barrier starts below $-\ell_c/\beta$ and changes by at
most $\|g\|\|h\||t|<|t|$, so it also stays feasible.

The proof of Lemma~\ref{lem:budget} now applies unchanged to the original
spectral potential and proves the same $L<L_*$ invariant.
The step-count proof uses \eqref{eq:screenedflat} in place of
\eqref{eq:flat} and still gives $O(SJ^2q^2)$ updates in a phase of size
$S$. Thus there is no new acceptance hypothesis.

Every refreshed warm row has $\chi_i>H_0-2$. Either
$|d_i|>(H_0-2)/2$ or $C_cr_i>(H_0-2)/2$. Combining
\eqref{eq:rowenergybound} and $\sum r_i^2\le\theta^2s$ gives
\begin{equation}\label{eq:warmcount}
 |\mathcal W|\le
 \frac{16384n+4C_c^2\theta^2s}{(H_0-2)^2}
 =O(n+s\log^2(2+n)).
\end{equation}
The local direction and progress arguments do not assume this row count;
the separate algebraic energy invariant supplies it for the cost analysis.

\subsection{The clock and its stronger path-length budget}\label{sec:clock}
Maintain a rational clock $\vartheta$, initially zero, adding $|t|$ for
every move $th$. It dominates path length because $\|h\|\le1/16$.
If a row was certified at time $\vartheta_i$ with score $\chi_i$ and
scale $r_i>0$, then while it remains cold,
\[
 \chi_i(\vartheta)\le\chi_i(\vartheta_i)
                        +r_i(\vartheta_i)(\vartheta-\vartheta_i).
\]
Moves change its tracked sum at speed at most $\|b_i\|\le r_i$;
cleanup preserves that sum and only lowers its scale. Store the rational
expiration time
\[
 \vartheta_i^{\rm exp}
 =\vartheta_i+\frac{H_0-\chi_i(\vartheta_i)}{r_i(\vartheta_i)}
\]
in a priority queue, with row-index tie breaking. An old key remains valid
through support deletion or scale decrease. Completed rows leave the
queue, and a row newly demoted from warm receives a new key.
The hysteresis rules make every cold renewal leave at least one unit of
score slack. If $r_i^{\rm ent}$ is its entry scale, successive due
inspections are at least $1/r_i^{\rm ent}$ apart in clock time, including
when an intervening warm period occurs.

Our update bound improves the clock estimate needed here. In a phase of
initial size $S$, exact $x^\T h=0$ and $\|h\|\ge1/32$ give
$\sum t^2\le1024S$. There are $O(SJ^2q^2)$ moves in the phase.
Cauchy--Schwarz, including short moves, therefore gives
\begin{equation}\label{eq:clockbudget}
 \sum_{\rm phase}|t|=O(SJq),\qquad
 \Lambda:=\int d\vartheta=O(nJq),\qquad
 \mathcal I:=\int s(\vartheta)\,d\vartheta=O(n^2Jq).
\end{equation}
Assign a drop in active dimension to the end of its move, so $s$ is
nonincreasing along the clock.

An inspection costs $O(s+\log(2+m_0))$. Its entry scale obeys
$r_i^{\rm ent}\le\theta\|a_i\|_2$, whence
$\sum_i r_i^{\rm ent}\le\theta\sqrt{m_0n}$.
For two consecutive due inspections at times $v'<v''$,
\[
 s(v'')\le r_i^{\rm ent}\int_{v'}^{v''}s(v)\,dv.
\]
The intervals for a fixed row are disjoint. Charging first inspections
separately and using \eqref{eq:clockbudget}, the total row-evaluation cost
of due inspections is
\begin{equation}\label{eq:inspectioncost}
 O(m_0n)+\widetilde O(\sqrt{m_0}\,n^{5/2}).
\end{equation}
The queue work is smaller: its total is at most
$\widetilde O(m_0+\sqrt{m_0}\,n^{3/2})$, plus support-event queue updates.
Finally
$\sqrt{m_0}\,n^{5/2}\le(m_0n^2+n^3)/2$,
so \eqref{eq:inspectioncost} is absorbed into
$\widetilde O(m_0n^2+n^3)$.

\subsection{Maintaining row energy without scanning all rows}
Row energy enters direction computation only through its gradient
$2(Qx_V+u)$ and the quadratic form $Q$. Computing them from the maintained
statistics costs $O(s^2)$ and requires no cold row evaluations.

Initialize $Q$ by classical multiplication in $O(m_0n^2)$ operations;
initially $u=0$. For a deliberate deletion of entry $a=b_{ij}$, let $b_i$
and $c_i$ denote the old row and offset. The exact updates are
\begin{align}\label{eq:statupdates}
 Q'&=Q-a(e_jb_i^\T+b_ie_j^\T)+a^2e_je_j^\T,\notag\\
 u'&=u+ax_jb_i-ac_ie_j-a^2x_je_j.
\end{align}
They cost $O(s)$. They follow by expanding the new retained row
$b_i-ae_j$ and new offset $c_i+ax_j$.
Freezing coordinate $j$, with $R=V\setminus\{j\}$, instead uses
\[
 Q'=Q_{R,R},\qquad u'=u_R+x_jQ_{R,j},\qquad
 c_i'=c_i+b_{ij}x_j.
\]
Use the old $Q$ in this formula, and do not also apply
\eqref{eq:statupdates} to the same freezing event.
Store active coordinates in a packed list with their original labels.
Swapping a frozen coordinate with the last active one and dropping that
row and column costs $O(s)$ for the statistics, without copying the full
principal submatrix. The corresponding row-record updates are included
in the support scan. Tie breaking continues to use original column labels.
The large-row diagonal $B$ is maintained at support events as well.

There are at most $n$ freezing events and $m_0n$ deliberate deletions.
At a freezing event one may simply scan all retained rows and recompute
their masses and scales, costing at most $O(m_0n)$ for that event.
Each deliberate deletion can be charged an $O(n)$ rescan and statistic
update. Scale searches and queue changes have logarithmic cost.
Thus all initialization, support maintenance, class changes, and forced
row inspections together cost $\widetilde O(m_0n^2)$.
After a move with no newly frozen coordinate, supports, masses, and scales
are unchanged, so no all-row cleanup scan is performed.

\subsection{Full screened procedure and operation count}
Use the following finite implementation of the core algorithm.
\begin{enumerate}
\item Scan the original matrix and keep rows with $\ell_1$ norm greater
than 66. Use the core's immediate returns when $m_0=0$ or $n\le512$.
Initialize the row state, $Q,u$, clock, and queue;
perform all initial cleanup with the statistic updates above.
\item At each anchor, refresh current warm rows and process all expired
cold certificates. Keep the large-row diagonal and statistics exact.
Form all spectral approximants from warm weights only, with
\eqref{eq:screeningprecision} and the stated warm error budget.
\item Impose the core constraints, \eqref{eq:chargegradient}, and the
computed high-$Q$ constraints. Form the negative basis and flatten just
coordinate vectors and warm medium gradients. Use the same dyadic step
schedule and full-or-boundary rule as the core algorithm, with the bound
$k>s/400$.
\item Advance the clock by the absolute step parameter. If coordinates
have frozen, perform cleanup and update $Q,u,B$ and the row records.
Apply the same stopping and phase-reset rules. The parameters
$\delta_0,\rho_c,\ell_c,C_c$ stay fixed throughout the run.
\item At $s\le512$, use the original conditional-expectation terminal
rounding on the surviving input matrix, retaining every original column.
\end{enumerate}
All classification tests are exact rational score comparisons. The only
transcendental approximations are the profile and exponential evaluations
already specified in the core. The discrepancy proof is unchanged:
the barriers, tracking identity, and terminal allowance use the same
constants, and the screened implementation preserves their invariants.
Discarded input rows satisfy the strict final bound because $66<C_*$.

Let $w_s$ be the number of warm rows at a state. Excluding support events
and due inspections, the field-operation cost per update is
$\widetilde O(w_s s^2+s^3)$, including the additional diagonalization of
$Q$. With fast primitives it is
$\widetilde O(w_s s^{\omega-1}+s^\omega)$.
By \eqref{eq:warmcount}, these are respectively
$\widetilde O(ns^2+s^3)$ and
$\widetilde O(ns^{\omega-1}+s^\omega)$.
A phase starting at $S$ has $O(SJ^2q^2)$ updates. Summing its geometrically
decreasing sizes gives total ordinary-update work
$\widetilde O(n^4)$ in field arithmetic and
$\widetilde O(n^{\omega+1})$ with fast primitives.
Add the one-time input scan, support maintenance, and
\eqref{eq:inspectioncost}. This proves
\[
 O(mn)+\widetilde O(n^4),\qquad
 O(mn)+\widetilde O(m_0n^2+n^{\omega+1}),
\]
respectively, and completes Theorem~\ref{thm:main}.

The term $m_0n^2$ is the remaining limitation for very tall inputs: the
proof explicitly pays for initializing and maintaining the retained-row
Gram statistics and support changes. We do not infer the stronger
universal bound $O(mn)+\widetilde O(n^{\omega+1})$ from screening alone.
For square input $m\le n$, that stronger expression does follow, because
$m_0n^2\le n^3$ is absorbed by $n^{\omega+1}$.

\begin{thebibliography}{9}
\bibitem{project} E. Akbas and S. Sra, project draft, version of
24 September 2026, especially the deterministic core and finite arithmetic
appendix; an extended version is
\nolinkurl{archive/algo-komlos-development.tex} in this repository; see
also the algorithmic appendix of arXiv:2609.27172v1.
\bibitem{core} E. Akbas and S. Sra, project note
\emph{Perron preconditioning and phase regularization for spectral
Koml\'os rounding}, 22 September 2026.
\bibitem{Li} X. Li, \emph{Fast Spectral Signing for Vector Balancing},
arXiv:2609.30044v1, 24 September 2026, Lemma 2.4 and Sections 4.3, 5, 9.
\url{https://arxiv.org/html/2609.30044v1}.
\end{thebibliography}
\end{document}
